\documentclass[10pt,a4paper]{amsart}
\usepackage[margin=19mm]{geometry}
\usepackage{amsmath,amssymb,mathtools}
\usepackage[scr=rsfs,cal=euler]{mathalfa}
\usepackage{xfrac}
\usepackage[unicode,hidelinks,bookmarks=false]{hyperref}
\newcommand{\ie}{i.e.\ }
\newcommand{\cf}{cf.\ }
\newcommand{\resp}{respectively\ }
\newcommand{\Subsection}{\subsection}
\providecommand{\nobreakdash}{\nobreak\hbox{-}}
\setlength{\emergencystretch}{2em}
\multlinegap0pt
\usepackage{amssymb}
\usepackage{mathtools}
\usepackage{xcolor}
\newtheorem{lemma}{Lemma}
\DeclareMathOperator{\arcosh}{arcosh}
\newcommand{\dd}{\,d}
\title{\bfseries The Balancing Theorem for Spanning Trees of Rectangular Grid Graphs}
\author{Jiechen Zhang}
\date{Source: arXiv:2605.23773v2 (CC BY 4.0)}
\begin{document}
\maketitle

\noindent\emph{Source and licence.} \bfseries The Balancing Theorem for Spanning Trees of Rectangular Grid Graphs. Version arXiv:2605.23773v2 (CC BY 4.0), distributed by arXiv under the Creative Commons Attribution 4.0 International licence, http://creativecommons.org/licenses/by/4.0/, as recorded in the arXiv licence metadata for this exact version and shown on the abstract page https://arxiv.org/abs/2605.23773v2. Full licence notice: \texttt{licenses/arxiv-2605.23773-CC-BY-4.0.txt}.\par\smallskip

\noindent\textit{Editorial scope.} Counted unit: the article's lemma lem:trap (source0.tex lines 401--417). The article's complete original proof of it is delivered with it (source0.tex lines 419--537, one \texttt{proof} environment of 119 lines). No mathematics is written, repaired or re-derived: the statement and the proof are the article's own lines, in the article's own notation, and no retained line is substituted or re-flowed.  No further source-local context is needed: every cross-reference of every retained line resolves inside the two delivered spans.  Nothing was omitted from any retained span, so the package carries no omission record.  No retained line contains a citation, so the wrapper carries no bibliography and no bibliography entry.  No retained line contains a figure environment, an image-inclusion command or drawing code, so no image is delivered and the package declares no asset dependency.  Editorial formatting only, in addition: an AMS \texttt{amsart} preamble that restates the article's own declarations and macros used by the retained lines, namely the environments \texttt{lemma} on separate counters, together with the standard packages those lines need.  The retained environments print in this document as Lemma 1 in order of appearance, whereas the article prints the same items with the numbering of its own full document.  The complete native source of the article as distributed in the arXiv e-print for this exact version is bundled unchanged as \texttt{sources/201234/source0.tex}.
\begin{lemma}\label{lem:trap}
For every $r\ge1$,
\begin{equation}\label{eq:trap-identity}
  T_r
  =
  I+
  \frac{1}{\pi r}
  \int_0^\pi
  \log\left(1-e^{-2r\gamma(\phi)}\right)\dd\phi.
\end{equation}
Consequently, if $1\le r\le s$, then
\begin{equation}\label{eq:C-trapezoidal}
  \frac{C_s}{s}-\frac{C_r}{r}
  \ge
  \frac{s-r}{rs}\alpha.
\end{equation}
\end{lemma}

\begin{proof}
We first prove \eqref{eq:trap-identity}. We use the standard identity, valid for $a\ge1$,
\[
  \frac1\pi\int_0^\pi\log(a-\cos\phi)\dd\phi
  =
  \log\frac{a+\sqrt{a^2-1}}{2}.
\]
With $a=2-\cos\theta$, this gives
\begin{equation}\label{eq:gamma-integral}
  \gamma(\theta)
  =
  \log2+
  \frac1\pi\int_0^\pi
  \log(2-\cos\theta-\cos\phi)\dd\phi.
\end{equation}

Fix $0<\phi\le\pi$ and put $q=e^{-\gamma(\phi)}$. Since
\[
  2-\cos\phi=\frac{q+q^{-1}}2,
\]
we have
\[
  2-\cos\theta-\cos\phi
  =
  \frac{1-2q\cos\theta+q^2}{2q}.
\]
Thus, with $L(\theta,\phi)=\log(2-\cos\theta-\cos\phi)$,
\[
  L(\theta,\phi)
  =
  \gamma(\phi)-\log2+
  \log(1-2q\cos\theta+q^2).
\]
For $0<q<1$,
\[
  \log(1-2q\cos\theta+q^2)
  =
  -2\sum_{m=1}^\infty\frac{q^m}{m}\cos(m\theta).
\]
The half-weighted trapezoidal filter is
\begin{equation}\label{eq:trap-filter}
  \frac1r\left(
  \frac{1+\cos(m\pi)}2+
  \sum_{j=1}^{r-1}\cos\frac{m\pi j}{r}
  \right)
  =
  \begin{cases}
  1,&2r\mid m,\\
  0,&2r\nmid m.
  \end{cases}
\end{equation}
Indeed, it is the real part of the periodic trapezoidal average
\[
  \frac1{2r}\sum_{\ell=0}^{2r-1}e^{i m\pi\ell/r}.
\]
Applying \eqref{eq:trap-filter} to the Fourier series of $\theta\mapsto L(\theta,\phi)$ gives
\begin{equation}\label{eq:L-average}
  \frac1r\left[
  \frac{L(0,\phi)+L(\pi,\phi)}2+
  \sum_{j=1}^{r-1}L(\pi j/r,\phi)
  \right]
  =
  \gamma(\phi)-\log2+
  \frac1r\log\left(1-e^{-2r\gamma(\phi)}\right).
\end{equation}

The preceding calculation is first made on $\phi\in[\varepsilon,\pi]$, where the Fourier series is uniformly convergent. Averaging \eqref{eq:gamma-integral} over the half-weighted $\theta$-grid and using \eqref{eq:L-average} gives the truncated identity on $[\varepsilon,\pi]$. It remains only to pass through the endpoint $\phi=0$. Since
\[
  \cos\phi=1-\frac{\phi^2}{2}+O(\phi^4),
  \qquad
  \gamma(\phi)=\arcosh\left(1+\frac{\phi^2}{2}+O(\phi^4)\right)
  =\phi+O(\phi^3),
\]
there is $\delta>0$ such that
\[
  \frac12\phi\le \gamma(\phi)\le 2\phi,
  \qquad 0<\phi\le\delta.
\]
Consequently, for fixed $r$,
\[
  \log\left(1-e^{-2r\gamma(\phi)}\right)
  =
  \log\phi+O_r(1),
  \qquad \phi\downarrow0.
\]
Equivalently, there is a constant $D_r>0$, depending only on $r$, such that
\[
  \left|\log\left(1-e^{-2r\gamma(\phi)}\right)\right|
  \le D_r(1+|\log\phi|),
  \qquad 0<\phi\le\delta,
\]
which is integrable at $0$.
On the left side, the only singular grid term is
\[
  L(0,\phi)=\log(1-\cos\phi)=2\log\phi+O(1);
\]
the remaining grid terms and $\gamma(\phi)-\log2$ are bounded near $0$. Hence both sides have at worst logarithmic endpoint singularities, with an integrable majorant independent of the cutoff $\varepsilon$. The contribution of $0<\phi<\varepsilon$ is $O_r(\varepsilon|\log\varepsilon|)$, and letting $\varepsilon\downarrow0$ gives \eqref{eq:trap-identity}.

It remains to prove monotonicity. For $z>0$ and $p>0$, set
\[
  h_p(z)=\frac1p\log(1-e^{-2pz}).
\]
Then
\[
  \frac{\partial}{\partial p}h_p(z)
  =
  -\frac1{p^2}\log(1-e^{-2pz})
  +
  \frac1p\frac{2z}{e^{2pz}-1}
  >0.
\]
Indeed, $0<1-e^{-2pz}<1$, so the logarithmic term is positive after the minus sign, and the second term is plainly positive. Since $\gamma(\phi)>0$ for $0<\phi\le\pi$, the already established identity \eqref{eq:trap-identity} gives, pointwise in $\phi$, $h_s(\gamma(\phi))\ge h_r(\gamma(\phi))$ whenever $s\ge r$. The endpoint bound above gives integrability, so no differentiation under the integral is needed. Therefore $T_s\ge T_r$, and
\[
  \frac{C_s}{s}+\frac{\alpha}{s}
  \ge
  \frac{C_r}{r}+\frac{\alpha}{r},
\]
which is equivalent to \eqref{eq:C-trapezoidal}.
\end{proof}

\end{document}
